\chapter{Automated Epistemic Soundness Audits: Table-to-Trace Parity \& Branchless AST Inspection}
\label{ch:epistemic_soundness}

The final architectural pillar of AMPS is the continuous enforcement of structural, mathematical, and computational integrity across both the codebase and documentation suite. Because generative AI swarms and reinforcement learning policies are intrinsically capable of generating ungrounded behaviors or exploiting simulator artifacts, a high-frequency financial microstructure platform requires rigorous, automated Continuous Integration and Continuous Deployment (CI/CD) verification gates. 

This chapter outlines the three automated CI/CD verification gates -- Branchless AST Inspection, Table-to-Trace Memory Parity, and Econometric Soundness -- and provides the canonical Data-Flow Matrix case studies that serve as the ground truth for engine validation.

\section{The Three Automated CI/CD Verification Gates}

To guarantee code integrity, execution performance, and econometric validity prior to any codebase merge, the AMPS build pipeline enforces three automated verification gates, executed via pre-commit hooks and local CI runners.

\begin{figure}[htbp]
\centering
\begin{tikzpicture}[
    scale=0.88, transform shape,
    box/.style={draw, rectangle, rounded corners, minimum width=3.4cm, minimum height=1.2cm, align=center},
    arr/.style={->, thick, >=stealth}
]
    \node[box, fill=blue!15] (gate1) {\textbf{Gate 1: Branchless AST}\\LLVM IR \& AST Inspection\\Zero \texttt{if/else} in Hot Paths};
    \node[box, fill=green!15] (gate2) [right=1.2cm of gate1] {\textbf{Gate 2: Memory Parity}\\Table-to-Trace Alignment\\Bitwise Arrow/Zarr Parity};
    \node[box, fill=purple!15] (gate3) [right=1.2cm of gate2] {\textbf{Gate 3: Econometric}\\NPIV Regularization\\KRR Adaptive Stability};

    \draw[arr] (gate1) -- (gate2);
    \draw[arr] (gate2) -- (gate3);
\end{tikzpicture}
\caption{The three automated CI/CD verification gates enforcing computational, data, and econometric invariants.}
\label{fig:three_ci_gates}
\end{figure}

\subsection{Gate 1: Branchless AST Auditing Gate}

As established in Chapter \ref{ch:data_oriented_design} and Chapter \ref{ch:distilled_swarm_scaling}, maintaining sub-0.2 microsecond latency across 1,000,000 agents requires eliminating branch misprediction penalties on CPU execution pipelines. 

Gate 1 executes a static code scanner that parses the Python Abstract Syntax Tree (AST) and the compiled LLVM Intermediate Representation (IR) generated by Numba for all \texttt{@njit} decorated functions in \texttt{app/engine/systems/}:

\begin{enumerate}
    \item \textbf{AST Tree Crawling:} The scanner traverses all AST nodes within matching kernels, order processing routines, and vectorized MLP evaluations. If it detects standard control flow branching (\texttt{ast.If}, \texttt{ast.IfExp}) within hot execution loops, the build immediately aborts with a syntax invariant violation.
    \item \textbf{Branchless SIMD Mask Enforcement:} Developers are required to express state transitions and conditional logic strictly through float masking:
    \begin{listing}[htbp]
    \caption{Branchless Float Masking Invariant.}
    \label{listing:branchless_float_masking}
    \begin{minted}{python}
# Permitted branchless SIMD transfer
fill_qty = desired_qty * (price <= best_ask_price) * is_alive_mask
    \end{minted}
    \end{listing}
    This guarantees that the Numba compiler generates branchless vector instructions (e.g., AVX2 \texttt{VANDPS}, \texttt{VMULPS}) rather than conditional branch jumps.
\end{enumerate}

\subsection{Gate 2: Table-to-Trace Memory Parity Gate}

As mandated by Rule 05-A (\texttt{05-data-flow-invariants.md}), every theoretical state transition described in the documentation must match the runtime numerical output of the simulation engine with zero tolerance for drift.

The automated test runner (\texttt{scripts/verify\_matrix\_trace\_parity.py}) executes integration suites against mock event streams:
\begin{enumerate}
    \item The test harness parses the Markdown Data-Flow Matrices in \texttt{docs/scientific\allowbreak\_model/} into structured DataFrames.
    \item The Numba JIT matching engine executes against a deterministic synthetic event feed (bypassing online LLM generation).
    \item The runner asserts point-by-point, bitwise numerical parity between the documented table values, the in-memory Apache Arrow buffers, and the disk-backed Zarr telemetry arrays.
\end{enumerate}

Any numerical divergence exceeding $10^{-5}$ halts the build, blocking regression commits.

\subsection{Gate 3: Automated Econometric Soundness Gate}

To ensure that the 1,000,000 Tier 2 student agents preserve realistic human-like bounded rationality without developing ungrounded high-frequency arbitrage exploits, Gate 3 executes automated econometric validation:

\begin{enumerate}
    \item \textbf{NPIV Causal Isolation:} Runs Non-Parametric Instrumental Variables regression with simulated HTTP 429 rate-limit instruments ($Z_{\text{HTTP429}}$) to verify that structural agent decisions respond causally to macroeconomic narrative shocks rather than endogenous spread fluctuations.
    \item \textbf{KRR Stability Auditing:} Evaluates Kernel Ridge Regression with Diagonal Adaptive Kernels across rolling scenario windows, asserting that dynamic kernel weights adapt to volatility jumps without numerical degeneracy.
\end{enumerate}

\section{Data-Flow Matrix Case Studies}

To ground the verification gates in physical reality, AMPS maintains exhaustive Data-Flow Matrix specifications across canonical stress regimes. These matrices define the exact empirical assertions validated by Gate 2.

\subsection{Case Study 1: The Liquidity Black Hole (Lehman Scenario)}

The most severe macroeconomic shock simulates the sudden insolvency of a Tier-1 prime broker or global clearing bank. The Macro Orchestrator broadcasts a breaking regulatory announcement. Algorithmic market makers immediately revise their toxic flow probability $\mu \to 1$. Under the Glosten-Milgrom model (Chapter \ref{ch:algorithmic_market_making}), spreads widen toward infinity as resting bids are withdrawn.

\begin{table}[htbp]
\centering
\small
\begin{tabular}{@{}llcc@{}}
\toprule
\textbf{Tick ($\Delta T$)} & \textbf{State Variable} & \textbf{Baseline Value} & \textbf{Expected Value} \\ \midrule
$T_0 + 10$ ms & Best Bid ($P_{\text{bid}}^{(1)}$) & $\$100.00$ & $\$99.50$ \\[0.2em]
$T_0 + 10$ ms & Best Ask ($P_{\text{ask}}^{(1)}$) & $\$100.02$ & $\$101.50$ \\[0.2em]
$T_0 + 50$ ms & L1 Bid Volume ($V_{\text{bid}}^{(1)}$) & $10,000$ shares & $2,500$ shares (75\% drain) \\[0.2em]
$T_0 + 150$ ms & Order Flow Imbalance (OFI) & $+0.05$ & $-0.92$ (Massive sell skew) \\[0.2em]
$T_0 + 200$ ms & Active Market Maker Count & $40$ agents & $4$ agents (90\% withdrawal) \\[0.2em]
$T_0 + 500$ ms & Bid-Ask Spread ($S$) & $\$0.02$ & $\$5.00$ ($250\times$ widening) \\ \bottomrule
\end{tabular}
\caption{The Lehman Scenario Data-Flow Matrix. Demonstrating liquidity voiding and defensive spread widening post-shock.}
\label{tab:dfm_lehman}
\end{table}

\subsection{Case Study 2: The Flash Rally (Short Squeeze)}

Market panic can manifest as an upward price explosion driven by forced short-covering. When unexpected positive regulatory news or sudden buyout interest is broadcast, highly leveraged short positions are liquidated. Aggressive market orders sweep resting Ask liquidity across all five Level-2 book levels, triggering exchange volatility halts.

\begin{table}[htbp]
\centering
\small
\begin{tabular}{@{}llcc@{}}
\toprule
\textbf{Tick ($\Delta T$)} & \textbf{State Variable} & \textbf{Baseline Value} & \textbf{Expected Value} \\ \midrule
$T_0 + 10$ ms & Order Flow Imbalance (OFI) & $+0.02$ & $+0.88$ (Massive buy skew) \\[0.2em]
$T_0 + 50$ ms & Ask Depth ($L_1 \dots L_5$) & $50,000$ shares & $0$ shares (Complete depletion) \\[0.2em]
$T_0 + 100$ ms & Best Ask ($P_{\text{ask}}^{(1)}$) & $\$50.00$ & $\$55.00$ (+10\% gap up) \\[0.2em]
$T_0 + 250$ ms & LULD Circuit Breaker & \texttt{False} & \texttt{True} (Trading halted) \\ \bottomrule
\end{tabular}
\caption{The Short Squeeze Data-Flow Matrix. Demonstrating rapid Ask book depletion and LULD circuit breaker triggering.}
\label{tab:dfm_short_squeeze}
\end{table}

\subsection{Case Study 3: The Fat-Finger Anomaly}

The third case study verifies that the simulation engine correctly isolates non-informational liquidity dislocations from genuine macroeconomic shocks. An institutional participant accidentally submits a sell order two orders of magnitude larger than intended (1,000,000 shares instead of 10,000), sweeping multiple bid price levels.

Because the Macro Orchestrator has broadcast no macroeconomic shock, informed trading probability $\mu$ remains low. Market makers temporarily widen spreads to absorb the imbalance, but rapidly re-quote when no secondary selling emerges, driving an immediate V-shaped price recovery.

\begin{table}[htbp]
\centering
\small
\begin{tabular}{@{}llcc@{}}
\toprule
\textbf{Tick ($\Delta T$)} & \textbf{State Variable} & \textbf{Baseline Value} & \textbf{Expected Value} \\ \midrule
$T_0 + 5$ ms & Executed Volume & $500$ shares & $1,000,000$ shares \\[0.2em]
$T_0 + 20$ ms & Best Bid ($P_{\text{bid}}^{(1)}$) & $\$100.00$ & $\$85.00$ (Instant dislocation) \\[0.2em]
$T_0 + 150$ ms & Market Maker Buy Intent & $2,000$ shares & $50,000$ shares (Stabilization) \\[0.2em]
$T_0 + 500$ ms & Best Bid ($P_{\text{bid}}^{(1)}$) & $\$85.00$ & $\$98.50$ (Mean reversion) \\ \bottomrule
\end{tabular}
\caption{The Fat-Finger Data-Flow Matrix. Verifying rapid V-shaped recovery following non-informational flow.}
\label{tab:dfm_fat_finger}
\end{table}

By enforcing bitwise parity against these three canonical scenario matrices in every CI/CD run, AMPS guarantees that the engine's microscopic order matching kernels faithfully reflect macroeconomic theory.
